% This is samplepaper.tex, a sample chapter demonstrating the
% LLNCS macro package for Springer Computer Science proceedings;
% Version 2.21 of 2022/01/12
%
\documentclass[runningheads]{llncs}
%
\usepackage[T1]{fontenc}
% T1 fonts will be used to generate the final print and online PDFs,
% so please use T1 fonts in your manuscript whenever possible.
% Other font encondings may result in incorrect characters.
%
\usepackage{graphicx}
% Used for displaying a sample figure. If possible, figure files should
% be included in EPS format.
%
% If you use the hyperref package, please uncomment the following two lines
% to display URLs in blue roman font according to Springer's eBook style:
%\usepackage{color}
\usepackage{kotex}
\usepackage{hyperref}
\usepackage{multirow}
\usepackage{listings}
\usepackage{xcolor}
%\renewcommand\UrlFont{\color{blue}\rmfamily}
%\urlstyle{rm}
%
\begin{document}
%
\title{Java에서의 Functional 프로그래밍의 현재 상태와 Haskell과의 비교}
%
%\titlerunning{Abbreviated paper title}
% If the paper title is too long for the running head, you can set
% an abbreviated paper title here
%
\author{Jinwoo Jeong}
%
% First names are abbreviated in the running head.
% If there are more than two authors, 'et al.' is used.
%
%
\maketitle              % typeset the header of the contribution
%
\begin{abstract}
Java는 Object-Oriented 프로그래밍 패러다임을 중심으로 돌아가는 언어입니다. 
본격적인 functional 프로그래밍과는 거리가 멀지만 Java 8에서 lambda 함수 등을 추가하는 등 노력을 하고 있습니다.
Vavr과 같이 FP 패러다임을 구현하려는 시도를 한 라이브러리도 있지만 Haskell과 같은 functional 프로그래밍 언어를 사용하는 편이 좋아 보입니다.

\keywords{Java \and Functional Programming \and Lambda Calculus \and Haskell}
\end{abstract}
%
%
%
\section{Introduction}
본 미니 연구는 학술적인 것이 아니며 Java에서의 functional 프로그래밍 구현을 알아보기 위한 것으로
reference로 Wikipedia\cite{ref_wikipedia}나 Baeldung\cite{ref_baeldung} 등을 폭 넓게 사용하였으니 참고 시 유의 바랍니다.

\noindent
\\
Java는 object를 중심으로 돌아가는 Object-Oriented 프로그래밍 언어입니다\cite{ref_wiki_object_oriented_programming}. 
그에 반해 functional 프로그래밍 패러다임은 function을 first-클래스로 취급하며 lambda calculus를 기반으로 하고 있습니다\cite{ref_wiki_functional_programming}.
모든 method가 class 안에서 생성되는 Java에서 functional interface를 통해 higher-order function을 구현하려는 등 여러 노력이 있었지만
Java를 이용해 본격적인 functional 프로그래밍을 하는 것은 쉽지 않아 보입니다. 
Vavr 라이브러리는 이러한 Java 환경 속에서 Object-Functional 프로그래밍\cite{ref_vavr_website}을 구현하기 위해 만들어졌습니다.
하지만 Java 컴파일러의 근본적인 차이로 인해, side-effect 없는 pure function을 구현하고자 한다면 Haskell과 같은 언어를 사용하는 것이 나아 보입니다.

\section{Background}

Java에서의 Functional 프로그래밍에 대해 분석하기에 앞서 Functional 프로그래밍이 무엇인지, 
그리고 그 기반이 되는 lambda calculus란 어떤지에 대해 개략적으로나마 알 필요가 있습니다.

\subsection{Lambda Calculus}

\noindent \\
Lambda calculus란 computation을 function의 abstraction 및 application을 통해 표현하는 formal system입니다.\cite{ref_wiki_lambda_calculus}. 
Untyped lambda calculus는 모든 종류의 computation을 표현 가능(i.e. Turing-complete)합니다. 
Typed lambda calculus의 경우 strongly normalizing(i.e. 모든 re-write 시퀀스가 normal form으로 reduce 됨)하는 경우 computation이 끝나게 되므로 모든 Turing-computation의 표현이 불가합니다 \cite{ref_wiki_typed_lambda_calculus}.

\noindent \\
Lambda calculus를 이해하기 위해 lambda term과 transformation 및 reduction에 대해 알아 봅시다. 아래 내용은 Wikipedia를 기반으로 합니다\cite{ref_wiki_lambda_calculus_definition}\cite{ref_wiki_lambda_calculus}\cite{ref_wiki_currying}.

\noindent \\
$\lambda x.M$을 예로 들어봅시다. 해당 expression에는 다음과 같은 lambda term이 있습니다:
\noindent \\
- $\lambda x.M$는 $\lambda$ abstraction이라고 부르며, 함수를 정의합니다. $\lambda$는 abstract operator라고 불립니다.
\noindent \\
- x를 variable이라 합니다. Free variable과 bounded variable로 구분할 수 있습니다.
\noindent \\
- M은 function의 body입니다. Input 없는 함수의 경우 단순히 M만으로 표현 될 수가 있습니다.

\noindent \\
- M N의 경우 application이라고 불립니다

\noindent \\
예를 들어 $\lambda x.x$는 Identity로 입력 값을 그대로 output하며  $\lambda x.x+1$은 입력값 + 1을 주는 function입니다. 
Variable은 free variable과 bound variable로 분류됩니다. 
Variable의 substitution 작업의 경우 $M[x:N]$의 형태로 표현되며, $\lambda y.M[x:=N]$의 경우 y가 N의 free variable의 set(i.e. FV(N))에 속할 경우 $\alpha$-renaming이 진행됩니다.

\noindent \\
$\alpha$-equivalence는 bound variable의 이름 변경이 되어도 같은 lambda term이라는 것을 표현 합니다.

\noindent \\
$\beta$-reduction은 $(\lambda x.t)s$ 형태의 application인 $\beta$-redex가 $t[x:=s]$로 reduce 되는 것입니다. 예시로:\\
$(\lambda x.x)s -> x[x:=s] -> s$가 있습니다.

\noindent \\
Currying은 여러 개의 argument를 받아들이는 function을 단일 argument를 받는 function의 sequence로 만드는 것입니다.

\subsection{Functional 프로그래밍}

Functional 프로그래밍은 lambda calulus를 구현하며 \cite{ref_wiki_functional_programming}, 
컴퓨터의 computation을 수학의 function을 evaluate하는 것과 같이 컴퓨터 프로그램을 작성하는 프로그래밍 패러다임입니다\cite{ref_baeldung_java_functional_programming}.

\noindent
\\
Functional 프로그래밍 언어의 대표적인 예로는 널리 쓰이는 언어 중 가장 functional하며 이후 Java 코드와 비교에 사용할 Haskell\cite{ref_haskell_website}이나 약간의 impurity를 더하지만 
Jane Street가 사용하는 언어로 유명한 OCaml\cite{ref_ocaml} 등이 있습니다.\\ 

\noindent
\\
Functional 프로그래밍의 주요 특징으로는 다음이 있습니다\cite{ref_wiki_functional_programming}\cite{ref_baeldung_java_functional_programming}: \\
- Function이 first-class로 취급됨 -> FP의 핵심 \\
- Hihger-order Function -> Function을 argument로 받음 OR function을 output함\cite{ref_wiki_higher_order_function}\\
- Pure function: Side effect가 없는 함수(메모리나 I/O 등). Haskell에서는 IO 작업에 Monad를 사용함.

\noindent
\\
언어 구현 상 및 FP로 개발시의 특징으로는 다음이 있습니다: \\
- Hindley-Miller type 시스템을 이용해 type 표현 -> Haskell 등\cite{ref_wiki_hindley_milner_type_system}\\
- Recustion: Loop의 경우 FP에서는 일반적으로 Recursion을 통해 이루어짐 \\

\section{Java에서의 Functional 프로그래밍}

\subsection{Java와 Lambda Expression}

$(λf. λg. λx. f (g x)) (+1) (*2) 5$와 같은 $\lambda$ expression을 예로 들어봅시다.

\noindent
\\
Haskell로는 다음과 같습니다.\\
\begin{lstlisting}[language=Haskell]
(\f -> \g -> \x -> f (g x)) (+1) (*2) 5 
\end{lstlisting}

\noindent
\\
Java로는 다음과 같습니다.
\begin{lstlisting}[language=Java]
import java.util.function.Function;
public class LambdaBetaReduction {
    static Function<Function<Integer,Integer>,
            Function<Function<Integer,Integer>,
                    Function<Integer,Integer>>> compose =
            f -> g -> x -> f.apply(g.apply(x));

    public static void main(String[] args) {
        int result = compose.apply(n -> n + 1)
                .apply(n -> n * 2)
                .apply(5);
    }
}    
\end{lstlisting}

\noindent
\\
Java8 이후 추가된 Function interface와 anonymous function을 활용하여야 합니다.

\subsection{Immutable 데이터와 pure function에 대해}

Referential transparency를 위해선 function이 pure해야 합니다. Function이 pure하기 위해서는 데이터가 immmutable 해야합니다.

\noindent
\\

\section{Java로 Haskell의 feature 따라하기}

\subsection{Java Optional과 Haskell Maybe}

\subsection{Side Effect에 대하여}

\section{결과}

Java8에서 Functional 인터페이스와 lambda expression이 추가된 이후, 자바에서 functional 프로그래밍을 하는 것이 한결 편해졌습니다.
하지만, Java의 경우 기본적으로 object-oriented 프로그래밍을 중심으로 한 언어로, 

\noindent
\\
또한, 

\section{후속 연구}

앞서까지 Java에서의 Functional 프로그래밍에 대하여 알아보았습니다. 

\section{사용된 툴 맟 코드}

해당 미니 연구를 포함한 GitHub repository는 \texttt{git@github.com:kosmostree/mini-researches.git}에서 clone 가능하며
 \href{https://github.com/kosmostree/mini-researches}{온라인}으로도 보실 수 있습니다.

\noindent
\\
본 연구에서 사용된 Java와 Haskell 툴은 다음과 같습니다.\\

\noindent Java: \\
- 컴파일러: javac 25.0.3 \\
- Java 버전: openjdk 25.0.3 2026-04-21 \\

\noindent Haskell: \\
- 컴파일러: The Glorious Glasgow Haskell Compilation System, version 9.10.3 \\

%
% ---- Bibliography ----
%
% BibTeX users should specify bibliography style 'splncs04'.
% References will then be sorted and formatted in the correct style.
%
% \bibliographystyle{splncs04}
% \bibliography{mybibliography}
%
\begin{thebibliography}{8}

\bibitem{ref_haskell_website}
Haskell 홈페이지, \url{https://www.haskell.org/}, last accessed 2026/07/13

\bibitem{ref_vavr_repo}
GitHub vavr-io/vavr 페이지, \url{https://github.com/vavr-io/vavr}, last accessed 2026/07/13

\bibitem{ref_vavr_website}
Vavr 홈페이지, \url{https://vavr.io/}, last accessed 2026/07/13

\bibitem{ref_wiki_lambda_calculus}
Wikipedia Lambda calculus 페이지, \url{https://en.wikipedia.org/wiki/Lambda_calculus}, last accessed 2026/07/14 

\bibitem{ref_wiki_lambda_calculus_definition}
Wikipedia Lambda calculus definition 페이지, \url{https://en.wikipedia.org/wiki/Lambda_calculus_definition}, last accessed 2026/07/14 

\bibitem{ref_wiki_functional_programming}
Wikipedia Functional programming 페이지, \protect\url{https://en.wikipedia.org/wiki/Functional_programming}, last accessed 2026/07/13

\bibitem{ref_wiki_object_oriented_programming}
Wikipedia Object-oriented Programming 페이지, \url{https://en.wikipedia.org/wiki/Object-oriented_programming}, last accessed 2026/07/13

\bibitem{ref_wiki_programming_paradigm}
Wikipedia Programming paradigm 페이지, \url{https://en.wikipedia.org/wiki/Programming_paradigm}, last accessed 2026/07/13

\bibitem{ref_baeldung_java_functional_programming}
Baeldung Functional Programming in Java 페이지,\url{https://www.baeldung.com/java-functional-programming}, last accessed 2026/07/14

\bibitem{ref_baeldung_java_lambda_expressions}
Baeldung  페이지, \url{https://www.baeldung.com/java-8-lambda-expressions-tips}, last accessed 2026/07/14

\bibitem{ref_baeldung_java_currying}
Baeldung Currying in Java 페이지, \url{https://www.baeldung.com/java-currying}, last accessed 2026/07/14

\bibitem{ref_ocaml}
OCaml 웹사이트, \url{https://ocaml.org/}, last accessed 2026/07/14

\bibitem{ref_wikipedia}
Wikipedia 웹사이트, \url{https://www.wikipedia.org/}, last accessed 2026/07/14

\bibitem{ref_baeldung}
Baeldung 웹사이트, \url{https://www.baeldung.com/}, last accessed 2026/07/14

\bibitem{ref_wiki_typed_lambda_calculus}
Wikipedia Typed lambda calculus 페이지, \url{https://en.wikipedia.org/wiki/Typed_lambda_calculus}, last accessed 2026/07/14

\bibitem{ref_wiki_currying}
Wikipedia Currying 페이지, \url{https://en.wikipedia.org/wiki/Currying}, last accessed 2026/07/15

\bibitem{ref_wiki_higher_order_function}
Wikipedia Higher-order function, \url{https://en.wikipedia.org/wiki/Higher-order_function}, last accessed 2026/07/15

\bibitem{ref_wiki_purely_functional_data_structure}
Wikipedia Purely functional data structure 페이지, \url{https://en.wikipedia.org/wiki/Purely_functional_data_structure}, last accessed 2026/07/15

\bibitem{ref_wiki_hindley_milner_type_system}
Wikipedia Hindley Milner type system 페이지, \url{https://en.wikipedia.org/wiki/Hindley%E2%80%93Milner_type_system}, last accessed 2026/07/15

\end{thebibliography}
\end{document}